% ============================================================================
\clearpage
\section{Regímenes y confianza medible}\label{sec:confianza}
% ============================================================================

\ficha{Errores de pronóstico ya madurados, factores filtrados, indicadores de calidad de datos y de
ajuste.}
{Detectar cambios de comportamiento, medir calibración, construir intervalos adaptativos y decidir si
abstenerse.}
{Probabilidad de cambio reciente, métricas de calibración, intervalos con nivel efectivo y estado de
confianza por dimensión.}
{Calibración en un tramo temporal separado; actualización solo con resultados madurados; abstención con
motivo registrado.}

\subsection{Detección de cambios de régimen}\label{sec:bocpd}

La detección bayesiana en línea de puntos de cambio \citep{adams2007} mantiene la distribución posterior
de la \emph{longitud de corrida} $r_t$, el número de sesiones desde el último cambio. Con una probabilidad
de riesgo $H$ por sesión y la predictiva $\pi_t^{(r)}=p(x_t\mid r_{t-1}=r,\ x^{(r)})$ calculada con los
datos de la corrida,
\begin{align}
P(r_t=r+1,\,x_{1:t}) &= P(r_{t-1}=r,\,x_{1:t-1})\;\pi_t^{(r)}\;(1-H),\label{eq:bocpd1}\\
P(r_t=0,\,x_{1:t}) &= \textstyle\sum_{r}P(r_{t-1}=r,\,x_{1:t-1})\;\pi_t^{(r)}\;H .\label{eq:bocpd2}
\end{align}
Con una previa normal--gamma inversa, la predictiva es una $t$ de Student y la recursión es exacta. Se
aplica sobre errores de pronóstico o factores, y se compara primero con un detector simple por umbral
(figuras~\ref{fig:umbral} a~\ref{fig:cambio_reciente}).

\grafica[0.6]{s5_umbral}{Errores estandarizados de pronóstico en un mundo sintético con un cambio de
varianza en la sesión 160, y alarmas de un detector por umbral (dos de las últimas cinco sesiones con
$|z|>2.5$).}{fig:umbral}

\grafica[0.6]{s5_corrida}{Posterior de la longitud de corrida de BOCPD en escala logarítmica. Tras el
cambio, la masa colapsa a corridas cortas y vuelve a crecer desde cero.}{fig:corrida}

\grafica[0.6]{s5_cambio_reciente}{Probabilidad posterior de que haya ocurrido un cambio en las últimas diez
sesiones. Es alta justo después del cambio verdadero y baja el resto del tiempo.}{fig:cambio_reciente}

\begin{noconcluir}
Una probabilidad alta de cambio de régimen dice que el comportamiento reciente difiere del aprendido; no
es una probabilidad de caída del precio. Su uso natural es degradar la confianza y ampliar la
incertidumbre, no generar señales direccionales.
\end{noconcluir}

\subsection{Cuatro dimensiones de confianza}\label{sec:dimensiones}

La salida del sistema distingue cuatro fuentes de duda que no deben mezclarse en un solo número
(\figref{fig:confianza}). Un alto número de contratos observados no equivale a muchas sesiones
independientes: cien strikes de un mismo día son una sola observación del estado del mercado
(\secref{sec:particiones}).

\begin{figure}[H]
\centering
\resizebox{\linewidth}{!}{% Cuatro dimensiones de confianza, cada una con indicadores y estados.
\begin{tikzpicture}[
  tarjeta/.style={draw=#1,fill=#1!5,rounded corners=3pt,line width=0.6pt,text width=6.9cm,
    minimum height=2.55cm,align=left,inner sep=6pt,font=\sffamily\scriptsize,text=tinta,anchor=north west},
  estado/.style={circle,fill=#1,minimum size=2.2mm,inner sep=0pt}]
\node[tarjeta=qazul] (c1) at (0,0) {\textbf{\small 1. Calidad de datos}\\[2pt]
  ¿Se puede reconstruir $\Q$ en el corte?\\[2pt]
  \textcolor{tinta2}{Spread relativo, strikes válidos por vencimiento, retraso de la fuente, snapshots
  faltantes, sincronía con el subyacente.}};
\node[tarjeta=qazul] (c2) at (7.45,0) {\textbf{\small 2. Sensibilidad del ajuste}\\[2pt]
  ¿Cuánto cambia $\Q$ dentro de lo compatible con los datos?\\[2pt]
  \textcolor{tinta2}{Ancho de bandas de cuantiles, discrepancia SSVI frente a ajuste convexo, restricciones
  activas, error de desamericanización.}};
\node[tarjeta=qnaranja] (c3) at (0,-2.85) {\textbf{\small 3. Incertidumbre predictiva}\\[2pt]
  ¿Qué tan dispersa y qué tan calibrada es la predicción $\P$?\\[2pt]
  \textcolor{tinta2}{Ancho de intervalos, cobertura y PIT recientes, probabilidad de cambio de régimen,
  desempeño del último tramo maduro.}};
\node[tarjeta=qaqua] (c4) at (7.45,-2.85) {\textbf{\small 4. Estabilidad de la decisión}\\[2pt]
  ¿Cambiaría la recomendación ante perturbaciones plausibles?\\[2pt]
  \textcolor{tinta2}{Dispersión de pesos óptimos remuestreados, distancia a la región casi óptima,
  sensibilidad a costos y a escenarios de estrés.}};
% Leyenda de estados (icono + etiqueta; el color nunca va solo)
\node[estado=bueno] (e1) at (0.25,-5.95) {};
\node[nota,anchor=west,text=tinta] at (e1.east) {\ \textbf{apto}: se usa sin restricciones};
\node[estado=alerta] (e2) at (5.0,-5.95) {};
\node[nota,anchor=west,text=tinta] at (e2.east) {\ \textbf{degradado}: solo reducciones o mantener};
\node[estado=critico] (e3) at (10.15,-5.95) {};
\node[nota,anchor=west,text=tinta] at (e3.east) {\ \textbf{bloqueado}: sin decisiones basadas en señal};
\end{tikzpicture}
}
\caption{Dimensiones de confianza y sus indicadores. Cada dimensión se reporta con un estado explícito y
su etiqueta; el color nunca comunica el estado por sí solo.}
\label{fig:confianza}
\end{figure}

\subsection{Calibración y reglas de puntuación}\label{sec:calibracion}

La calibración de probabilidades se ajusta en un tramo temporal separado del de entrenamiento y del de
prueba. Las reglas de puntuación son propias \citep{gneiting2007}:
\begin{align}
\text{Brier} &= \tfrac1N\textstyle\sum_i (p_i-o_i)^2 \quad\text{para eventos binarios,}\label{eq:brier}\\
\mathrm{CRPS}(F,y) &= \int_{-\infty}^{\infty}\bigl(F(z)-\mathbb 1\{y\le z\}\bigr)^2dz
  = 2\int_0^1\rho_\tau\bigl(y-F^{-1}(\tau)\bigr)\,d\tau ,\label{eq:crps}
\end{align}
y la pérdida cuantílica $\rho_\tau$ para percentiles individuales. La segunda igualdad de \eqref{eq:crps}
conecta directamente los pronósticos cuantílicos con el CRPS \citep{gneiting2011}. Las puntuaciones se
acompañan de curvas de calibración, cobertura y ancho de intervalos (figuras~\ref{fig:fiabilidad}
y~\ref{fig:pit}).

\grafica[0.58]{s5_fiabilidad}{Diagrama de fiabilidad de un pronóstico de evento en el mundo sintético.
El modelo sobreconfiado empuja las probabilidades hacia los extremos: su curva es más plana que la
diagonal y su Brier, peor.}{fig:fiabilidad}

\grafica[0.58]{s5_pit}{Histogramas PIT de pronósticos de distribución. Un pronóstico calibrado da valores
$F(y)$ uniformes; intervalos demasiado estrechos producen forma de U y demasiado anchos, una
joroba.}{fig:pit}

\subsection{Conformal adaptativo}\label{sec:conformal}

Como envoltura de intervalos se prueba la inferencia conformal adaptativa \citep{gibbs2021}. Si
$\mathrm{err}_t=\mathbb 1\{y_t\notin\widehat C_t(\alpha_t)\}$, el nivel efectivo se actualiza como
\begin{equation}
\alpha_{t+1}=\alpha_t+\gamma\,(\alpha-\mathrm{err}_t),
\qquad
\Bigl|\tfrac1N{\textstyle\sum_{t}}\mathrm{err}_t-\alpha\Bigr|\le\frac{\max\{\alpha_1,1-\alpha_1\}+h\,\gamma}{\gamma\,N},
\label{eq:aci}
\end{equation}
donde $N$ es el número de intervalos evaluados y $h$ el horizonte. Con horizonte $h$, el error de la
etiqueta que vence en $t$ solo se conoce en $t$: la actualización usa exclusivamente errores madurados
(\figref{fig:alineacion}), y los $h-1$ intervalos ya emitidos con un nivel anterior explican el término
$h\gamma$; con $h=1$ es la cota de \citet{gibbs2021}. El umbral de cada intervalo es el
$\lceil(n+1)(1-\alpha_t)\rceil$-ésimo puntaje de calibración.

La cota vale para \emph{cualquier} secuencia, sin supuestos sobre la distribución, pero exige dos
convenciones: con $\alpha_t\ge1$ el conjunto es vacío (error seguro) y con $\alpha_t\le0$ es la recta
completa. Así $\alpha_t$ queda en $[-h\gamma(1-\alpha),\,1+h\gamma\alpha]$. Sin la primera convención la
cota falla: si los puntajes empatan en cero, por ejemplo con retornos nulos, un intervalo de ancho cero
siempre cubre, $\alpha_t$ crece sin límite y la tasa de error se queda en cero. La primera versión del
núcleo tenía ese defecto. Ahora las pruebas comprueban la cota en cada prefijo con secuencias límite
(empates, puntajes siempre crecientes, alternancias, un adversario que ve el intervalo y huecos en los
datos), además de datos aleatorios, que por sí solos no lo detectaban. Parte de la cobertura puede venir
de intervalos vacíos o infinitos, que no informan: la tasa agregada no basta para juzgar los intervalos.

\grafica[0.62]{s5_conformal_intervalos}{Intervalos del 90\% ante un cambio de volatilidad en la sesión 500
y otro en la 750 (datos sintéticos). El método de nivel fijo tarda en ajustar su ventana; el adaptativo
reacciona antes.}{fig:conformal_intervalos}

\grafica[0.62]{s5_conformal_cobertura}{Cobertura móvil de 50 sesiones. Ambos métodos terminan cerca del
90\% en el promedio de largo plazo, pero el de nivel fijo pasa muchas sesiones por debajo del 75\%: la
garantía agregada oculta fallas locales.}{fig:conformal_cobertura}

\grafica[0.58]{s5_conformal_alfa}{Trayectoria del nivel efectivo $\alpha_t$: baja cuando los errores se
acumulan y sube cuando sobra cobertura.}{fig:conformal_alfa}

\begin{noconcluir}
El control de cobertura a largo plazo no implica que mañana, o en cada régimen, exista exactamente la
cobertura nominal; tampoco convierte el intervalo en una probabilidad de rentabilidad.
\end{noconcluir}

\subsection{Abstención}\label{sec:abstencion}

Una entrada nueva requiere tres cosas a la vez: datos aptos, estabilidad suficiente del ajuste y evidencia
predictiva relevante (\figref{fig:abstencion}). Si falta cualquiera, se bloquean los incrementos basados
en señal y se comunica el problema. Las reducciones exigidas por límites de riesgo definidos de antemano se
evalúan aparte, con datos adecuados.

\begin{figure}[H]
\centering
\resizebox{0.92\linewidth}{!}{% Regla de abstención: los incrementos exigen tres condiciones; las reducciones por límites van aparte.
\begin{tikzpicture}[
  pregunta/.style={diamond,aspect=2.6,draw=qazul,fill=qazul!7,line width=0.6pt,align=center,
    font=\sffamily\scriptsize,text=tinta,inner sep=1.2pt,text width=2.7cm},
  bloque/.style={caja=critico,text width=3.6cm},
  permite/.style={caja=bueno,text width=3.6cm},
  node distance=6mm and 9mm]
\node[caja=tenue,text width=3.4cm] (ini) {Nueva sesión: snapshot, ajuste y pronósticos};
\node[pregunta,below=of ini] (d1) {¿Datos aptos?\\ cobertura, spreads, retraso};
\node[pregunta,below=of d1] (d2) {¿Ajuste estable?\\ bandas, contraste de métodos};
\node[pregunta,below=of d2] (d3) {¿Evidencia predictiva\\ relevante y calibrada?};
\node[permite,below=of d3] (ok) {Se permiten incrementos basados en señal, dentro de límites y fuera de la zona de no operación};
\node[bloque,right=of d1] (b1) {Bloquear incrementos basados en señal y comunicar el problema};
\node[caja=alerta,text width=3.6cm,right=of d3] (b3) {Sin nuevas entradas: mantener o reducir; registrar el motivo};
\draw[flecha] (ini) -- (d1);
\draw[flecha] (d1) -- node[nota,right]{sí} (d2);
\draw[flecha] (d2) -- node[nota,right]{sí} (d3);
\draw[flecha] (d3) -- node[nota,right]{sí} (ok);
\draw[flecha] (d1) -- node[nota,above]{no} (b1);
\draw[flecha] (d2.east) -| node[nota,pos=0.25,above]{no} (b1.south);
\draw[flecha] (d3) -- node[nota,above]{no} (b3);
% Rama independiente de límites de riesgo
\node[caja=qaqua,text width=3.5cm,left=14mm of d2] (lim) {Límites de riesgo definidos de antemano (concentración, CVaR, pérdidas)};
\node[caja=qaqua,text width=3.5cm,below=of lim] (red) {Evaluar reducciones con datos adecuados, aunque el modelo se abstenga};
\draw[flecha] (lim) -- (red);
\node[nota,text width=3.6cm,below=1mm of red] {abstenerse del modelo no significa ignorar una exposición existente};
\end{tikzpicture}
}
\caption{Regla de abstención. La rama de la izquierda no depende del modelo: abstenerse de la señal no
significa ignorar una exposición existente.}
\label{fig:abstencion}
\end{figure}

\begin{criterio}
Calibración aceptable en el tramo reservado, cobertura estable por subperiodos y regla de abstención
ejercitada en sesiones con datos degradados sin producir decisiones basadas en señal.
\end{criterio}
